% ============================================================
% INTRO TEXT: CTI + AI/LLMs MILESTONES (DROP-IN)
% ============================================================

\subsection{Cyber Threat Intelligence and the Case for Domain-Specialized LLMs}
Cyber Threat Intelligence (CTI) sits at the intersection of rapidly evolving adversary behavior, fragmented technical evidence, and high-stakes operational decision-making. Enterprise security operations must continuously ingest unstructured text (vendor reports, malware writeups, incident postmortems, threat bulletins), correlate it with semi-structured telemetry (SIEM/EDR alerts, logs), and translate it into actionable determinations (e.g., ``what is happening, how severe is it, who is likely responsible, what should we do next''). In practice, this translation is bottlenecked by analyst time and by the high volume of noisy detections: alert fatigue is widely recognized as a core driver of delayed response and missed true positives, with false-positive-heavy queues increasing response times and overall risk \cite{ban2021combat,gelman2023teq}. Recent operational ML research illustrates why this matters at scale: on real-world SOC data, prioritization/triage systems can suppress large fractions of false positives while maintaining high detection rates and reducing incident queue time and resolution time \cite{gelman2023teq}---exactly the types of improvements organizations track through SOC KPIs such as precision/accuracy, false-alert reduction, and mean-time-to-respond/repair (MTTR).

A major milestone enabling CTI automation has been the maturation of shared \emph{representations} and \emph{taxonomies}. Knowledge bases like MITRE ATT\&CK provide a community-grounded ontology of adversary tactics and techniques that supports consistent reporting, threat-informed defense, and structured mapping from observations to known TTPs \cite{strom2020attack}. In parallel, standard interchange formats and protocols such as STIX 2.1 and TAXII 2.1 have enabled scalable sharing of threat intelligence objects across organizations and tools \cite{oasis2021stix,oasis2021taxii}. Vulnerability intelligence has similarly standardized around identifiers and structured repositories (e.g., CVE identifiers and the NVD), supporting automation for vulnerability management and measurement \cite{cve2026cve,nist2026nvd}. However, these standards also increase the cognitive and operational surface area: CTI workflows require accurate grounding to evolving identifiers (CVE/CWE), consistent mapping to ATT\&CK technique/mitigation IDs, and faithful extraction of entities and relations from heterogeneous narrative text---all of which are brittle under hallucination and minor semantic errors.

The LLM era introduced a compelling hypothesis: that large generative models can serve as a universal interface layer between unstructured CTI narratives and structured enterprise workflows (triage, enrichment, ATT\&CK mapping, mitigation drafting). Early CTI-focused LLM research formalized this into datasets and evaluations. SEvenLLM contributed bilingual CTI instruction data and an evaluation framework targeting incident analysis and response-style tasks \cite{ji2024sevenllm}. CTIBench provided one of the first CTI-specific benchmark suites aimed at measuring applied CTI knowledge and task competence, explicitly highlighting reliability and hallucination concerns that limit direct deployment of general-purpose LLMs in CTI pipelines \cite{alam2024ctibench}. Building on the need for more realistic and continuously relevant evaluation, AthenaBench extended CTIBench with improved dataset construction, refined metrics, and a new risk-mitigation task, while emphasizing that CTI benchmarks can become stale unless refreshed as the threat landscape changes \cite{alam2025athenabench}. AthenaBench also underscores an important empirical pattern: even strong proprietary LLMs can remain subpar on reasoning-intensive CTI tasks such as threat actor attribution and risk mitigation, motivating training regimes and model designs explicitly tailored to CTI workflows \cite{alam2025athenabench}.

These benchmarks help clarify why CTI is not a ``prompting problem'' alone, and why enterprise-grade CTI requires \emph{specialized} language models. First, CTI language is dense with domain-specific entities, abbreviations, and implicit causal structure (exploits $\rightarrow$ weaknesses $\rightarrow$ tactics $\rightarrow$ mitigations), and small factual slips can drive incorrect operational decisions. Second, SOC and CTI tooling increasingly demands \emph{structured} outputs (JSON extractions, ATT\&CK IDs, prioritized hypotheses, mitigation playbooks) that must be both syntactically correct and semantically grounded. Third, at scale, cost, latency, privacy, and on-prem deployment constraints become first-order concerns for many security organizations, increasing the appeal of smaller, domain-expert models rather than relying exclusively on frontier closed models.

Recent cybersecurity-specialized IFT research exemplifies this trajectory. CyberPal.AI introduces an expert-driven cybersecurity instruction dataset (SecKnowledge) and demonstrates that instruction fine-tuning can produce security-specialized models that substantially outperform their base counterparts on security tasks \cite{levi2024cyberpal}. The follow-on work, \emph{Toward Cybersecurity-Expert Small Language Models}, advances this direction by training CyberPal 2.0 as a family of cybersecurity-expert small language models (SLMs) and by emphasizing higher-fidelity, task-grounded reasoning traces for security workflows \cite{levi2025cyberpal2}. Taken together with CTI benchmarking work (CTIBench/AthenaBench), these milestones support a coherent research agenda: (i) define realistic CTI tasks and continuously updated evaluations, (ii) generate domain-grounded instruction corpora aligned to CTI/SOC workflows, and (iii) train deployable specialized models that can measurably improve operational outcomes---higher triage precision and consistency, fewer false alerts escalated to humans, and faster investigation/containment cycles (lower MTTR) \cite{alam2025athenabench,gelman2023teq,levi2025cyberpal2}.



% ============================================================
% Revised Abstract text (DROP-IN)
% ============================================================

Security analysts routinely analyze noisy, unstructured security artifacts into standardized, automation-ready representations for further analysis and/or processing. Although large language models (LLMs) have shown great promise for these tasks, there remains much room for improvement when it comes to reliably reproducing high quality results across a range of standard security operations and threat intelligence use cases. Improving the performance of LLMs for this domain  have largely relied on supervised fine-tuning (SFT) methods to date. However, in many cases open source, community-maintained resources (OSINT) define canonical identifiers and schemas that enable deterministic verification of model outputs, outputs that can be leverages to strengthen the performance of LLMs for a range of Cybersecurity tasks. In this study, we introduce \textit{Minerva} which leverages this structured, publicly available OSINT data, and use it to train a variety of LLMs using reinforcement learning with verifiable rewards (RLVR), an approach which improves the capability of LLMs across a range of core Cybersecurity and CTI tasks. \textit{Minerva} uses a unified dataset and training pipeline spanning multiple Cybersecurity subtasks, each paired with task-specific verifiers that score structured outputs and identifier predictions. In addition, in order to address reward sparsity during implementation, we propose a lightweight self-training mechanism that generates additional verified trajectories and distills them back into the model. Our experiments across a range of LLM families illustrate the strength of this method of training over more standard SFT approaches across multiple benchmarks.
